\subsection{Adjunctions}
\label{subsec:prelim-adjunctions}

\gutentag{0008}%
Let $\Adj$ be the free homotopy coherent adjunction of Riehl--Verity \cite[\S
3.1]{riehlverity2016adjunctions}, a simplicial category whose hom-simplicial
sets are quasi-categories, regarded as an $(\infty,2)$-category
\cite{haugseng2015rectification}. It has objects $-$ and $+$,
$1$-morphisms $l\colon - \to +$, $r\colon + \to -$ and $2$-morphisms
$\eta\colon \mathrm{id}_{-} \Rightarrow rl$,
$\varepsilon\colon lr \Rightarrow \mathrm{id}_{+}$
satisfying the triangle identities. By
\cite[Cor.~3.3.5]{riehlverity2016adjunctions} it is isomorphic to the
Schanuel--Street free adjunction \cite{schanuelstreet1986free},
so its hom-simplicial sets are nerves of categories. Let
$\ell,\rho\colon \Delta^{1} \to \Adj$
select $l$ and $r$.

\begin{defi}
\label{def:adjunction}
\gutentag{0009}%
An \emph{adjunction} in an $(\infty,2)$-category $\A$ is a functor $A\colon \Adj
\to \A$. We write it as $A(l) \dashv A(r)$,
with \emph{unit} $A(\eta)$ and \emph{counit} $A(\varepsilon)$.
A $1$-morphism of $\A$ is a \emph{left adjoint}, resp.\ a \emph{right adjoint},
if it lies in the image of
\begin{equation*}
\ell^{*},\ \rho^{*}\colon \mapp_{\twoCat}(\Adj,\A) \longrightarrow
\mapp_{\twoCat}(\Delta^{1},\A) \ ,
\end{equation*}
respectively.
\end{defi}

\gutentag{000B}%
Functors preserve adjunctions. In particular, for $t \in \A$ the functors
$\Hom_{\A}(t,-)\colon \A \to \Cat$
and
$\Hom_{\A}(-,t)\colon \A\opcat \to \Cat$
\cite{hinich2020yoneda,heine2023enriched} take an adjunction $l \dashv r$ to
adjunctions of $\infty$-categories
\begin{equation}
\label{eq:hom-adjunctions}
l_{\natural} \dashv r_{\natural}\colon \Hom_{\A}(t,x) \rightleftarrows
\Hom_{\A}(t,y) \ , \qquad r^{\natural} \dashv l^{\natural}\colon \Hom_{\A}(x,t)
\rightleftarrows \Hom_{\A}(y,t) \ ,
\end{equation}
where for the second we use the equivalence
$\Adj \simeq \Adj\opcat$,
$l \mapsto r\opcat$, $r \mapsto l\opcat$,
which fixes the objects and preserves the unit and counit
\cite[Rem.~4.3]{haugseng2021laxmonads}.

\gutentag{00FD}%
We will use the following lifting theorem, proved model-independently by
Abell\'an (and by Riehl--Verity
\cite[Thms.~4.3.11, 4.4.18]{riehlverity2016adjunctions} for
quasi-categorically enriched categories).

\begin{thm}[{\cite[Thm.~5.1]{abellan2026bifibrations}}]
\label{thm:adj-lifting}
\gutentag{000A}%
Let $\A$ be an $(\infty,2)$-category.
\begin{enumerate}
\item The maps $\ell^{*}$ and $\rho^{*}$ are monomorphisms, so being a left or a
right adjoint is a property of a $1$-morphism.
\item Consider tuples
$(l,r,\eta,\varepsilon,\omega)$
of $1$-morphisms
$l\colon x \to y$
and
$r\colon y \to x$,
$2$-morphisms
$\eta\colon \mathrm{id}_{x} \Rightarrow rl$
and
$\varepsilon\colon lr \Rightarrow \mathrm{id}_{y}$
satisfying the triangle identities
$\varepsilon l \circ l\eta \simeq \mathrm{id}_{l}$
and
$r\varepsilon \circ \eta r \simeq \mathrm{id}_{r}$,
and an equivalence
$\omega\colon r\varepsilon \circ \eta r \simeq \mathrm{id}_{r}$.
The map
$A \mapsto (A(l),A(r),A(\eta),A(\varepsilon),\omega_{A})$,
where $\omega_{A}$ is the image of the triangle identity
$r\varepsilon \circ \eta r = \mathrm{id}_{r}$
of $\Adj$, is an equivalence from
$\mapp_{\twoCat}(\Adj,\A)$
to the space of such tuples. In particular, every such tuple extends to an
adjunction $l \dashv r$ with unit $\eta$ and counit $\varepsilon$, uniquely up
to a contractible space of choices.
\item For
$l\colon x \to y$,
the map
$A \mapsto (A(r),A(\eta))$
from the space of adjunctions $A$ with $A(l) = l$ to the space of pairs
$(r,\eta)$ with $\eta$ the unit of an adjunction $l \dashv r$ is an
equivalence, and both spaces are empty or contractible. In particular, for
given $r$ and $\eta$, the space of adjunctions $l \dashv r$ with unit $\eta$ is
empty or contractible.
\end{enumerate}
\end{thm}

\begin{proof}
\gutentag{00FE}%
(1) By \cite[Thm.~5.1]{abellan2026bifibrations}, $\ell^{*}$ is the map on
underlying spaces of an equivalence of $\func(\Adj,\A)$ with a locally full
sub-$(\infty,2)$-category of $\func(\Delta^{1},\A)$, whose objects are the $l$
admitting $r$, $\eta$, $\varepsilon$ as in (2). Locally full inclusions are
monomorphisms, and so are their maps on underlying spaces. For $\rho^{*}$, apply
this to $\A\opcat$, using
$\Adj \simeq \Adj\opcat$.

(2) and (3) Both sides of (2) lie over
$\mapp(\Delta^{1},\A)$
via $l$, so it suffices to compare the fibres over each
$l\colon x \to y$.
If $l$ is not a left adjoint, both are empty, as the objects of the image in
(1) are the $l$ admitting a tuple. Otherwise, by (1), the fibre of the left-hand
side is contractible; let $A$ be in it, with $r_{0} = A(r)$,
$\eta_{0} = A(\eta)$
and
$\varepsilon_{0} = A(\varepsilon)$.
The transposition
$\mapp(r_{0},r') \simeq \mapp(\mathrm{id}_{x},r'l)$,
$\theta \mapsto \theta l \circ \eta_{0}$,
of
$r_{0}^{\natural} \dashv l^{\natural}$
\eqref{eq:hom-adjunctions} is natural in
$r' \in \Hom_{\A}(y,x)$.
By the triangle identities, as for $2$-categories, it identifies the invertible
$\theta$ with the $\eta'$ admitting an $\varepsilon'$ satisfying the triangle
identities.
\todo{2-categorical input: check that uniqueness of adjoints works verbatim.}
These $\eta'$ include the units of adjunctions
$l \dashv r'$. So the space of pairs $(r',\eta')$ with $\eta'$ admitting such
an $\varepsilon'$ is the space of pairs
$(r',\theta\colon r_{0} \simeq r')$,
which is contractible; this proves (3).
For the fibre of the right-hand side of (2), it remains to see that for
$(r',\eta') = (r_{0},\eta_{0})$
the pairs
$(\varepsilon',\omega)$
form a contractible space. The transposition
$\mapp(lr_{0},\mathrm{id}_{y}) \simeq \mapp(r_{0},r_{0})$,
$\varepsilon' \mapsto r_{0}\varepsilon' \circ \eta_{0}r_{0}$,
of
$l_{\natural} \dashv (r_{0})_{\natural}$
\eqref{eq:hom-adjunctions} is an equivalence, so the pairs
$(\varepsilon',\omega\colon r_{0}\varepsilon' \circ \eta_{0}r_{0} \simeq
\mathrm{id}_{r_{0}})$
form its fibre over $\mathrm{id}_{r_{0}}$, which is contractible. It contains
$(\varepsilon_{0},\omega_{A})$,
so every such $\varepsilon'$ is equivalent to
$\varepsilon_{0}$ and also satisfies the other triangle identity.
\end{proof}

\begin{defi}
\label{def:reflection}
\gutentag{000C}%
An adjunction $l \dashv r$ in $\A$ is a \emph{reflection} if its counit
$\varepsilon$ is invertible, and a \emph{coreflection} if its unit $\eta$ is
invertible.
The left adjoint in a reflection is called a \emph{localization},
and the right adjoint in a coreflection a \emph{colocalization}.
A $1$-morphism is \emph{reflective} if it is the right adjoint of a reflection,
and \emph{coreflective} if it is the left adjoint of a coreflection.
\end{defi}

\gutentag{000D}%
By \cref{thm:adj-lifting}(1),
being reflective is a property of a $1$-morphism $r$,
independent of the choice of a left adjoint,
and similarly for being coreflective.
As $r_{\natural}$ and $l_{\natural}$ on $\Hom_{\A}(t,-)$ form an adjunction with
counit
$\varepsilon_{\natural}$
and unit $\eta_{\natural}$ \eqref{eq:hom-adjunctions},
a right adjoint $r$ is reflective if and only if every $r_{\natural}$ is fully
faithful, and a left adjoint $l$ is coreflective if and only if every
$l_{\natural}$ is fully faithful.

\gutentag{00FF}%
Now we turn to an alternative formulation of adjunctions.
The idea is to reformulate the equivalence
\begin{equation*}
\Hom _{C} (Lx, y) \simeq \Hom _{D} (x, Ry)
\end{equation*}
from adjunctions of functors between $ \infty $-categories,
in a general $ (\infty, 2) $-category $ \E $.
We will make the technical assumption that $ \E $ admits lax pullbacks
for the remainder of this subsection.

\gutentag{001A}%
We will use the category of spans to organize the combinatorics.
Let
$\Lambda := \{0 \leftarrow 2 \rightarrow 1\}$
be the walking span and
$\Span(\E) := \func(\Lambda,\E)$
the $(\infty,2)$-category of spans in $\E$. We write a span as
$S = (X \xleftarrow{p} E \xrightarrow{q} Y)$,
a span from $X$ to $Y$, and $\Span_{\E}(X,Y)$ for the fibre of
$\Span(\E) \to \E \times \E$,
evaluation at $0$ and $1$, over $(X,Y)$.
A $1$-morphism $S \to S'$ of $\Span(\E)$ over
$(a,b) \in \Hom(X,X') \times \Hom(Y,Y')$
is a \emph{map of spans}: a $1$-morphism $c\colon E \to E'$ with equivalences
$p'c \simeq ap$ and $q'c \simeq bq$. By
\cite[Prop.~4.1.2]{abellangagnahaugseng2024straightening},
applied to strong transformations,
\begin{equation*}
\Hom_{\Span(\E)}(S,S') \simeq \Hom(X,X') \times_{\Hom(E,X')} \Hom(E,E')
\times_{\Hom(E,Y')} \Hom(Y,Y') \ ,
\end{equation*}
and the hom-$\infty$-categories of $\Span_{\E}(X,Y)$ are the fibres over
$(\mathrm{id}_{X},\mathrm{id}_{Y})$.
The underlying $\infty$-category of $\Span_{\E}(X,Y)$ is the fibre of
$\func(\Lambda,\iota_{1}\E)$
over $(X,Y)$, so a $1$-morphism of $\Span_{\E}(X,Y)$ is an equivalence if and
only if it is one in $\E$.

\gutentag{0019}%
\textbf{Notation.} For $1$-morphisms $f\colon X \to Y$ and $g\colon Z \to Y$ of
an $(\infty,2)$-category $\E$ we write
$(f \downarrow g) := X \ratimes_{Y} Z$
for their lax pullback, with projections $p$, $q$ and $2$-morphism
$\alpha\colon fp \Rightarrow gq$,
and
$(u,v,\varphi)\colon T \to (f \downarrow g)$
for the $1$-morphism given by $u\colon T \to X$, $v\colon T \to Z$ and
$\varphi\colon fu \Rightarrow gv$.
We consider $ (f\downarrow g) $ as the span
$ X \xleftarrow{p} (f\downarrow g) \xrightarrow{q} Z $
from $ X $ to $ Z $.
By the universal property of lax pullbacks,
$\Hom(T,(f \downarrow g)) \simeq (f_{\natural} \downarrow g_{\natural})$,
naturally in $T$.
For $\E = \Cat$, $(f \downarrow g)$ is the comma $\infty$-category.
For $f\colon X \to Y$ and $g\colon Y \to X$ we write
$(f \downarrow Y) := (f \downarrow \mathrm{id}_{Y})$
and
$(X \downarrow g) := (\mathrm{id}_{X} \downarrow g)$;
both are spans from $X$ to $Y$. We write $p$, $q$,
$\alpha\colon fp \Rightarrow q$
for the structure of $(f \downarrow Y)$, and $\bar p$, $\bar q$,
$\beta\colon \bar p \Rightarrow g\bar q$
for that of $(X \downarrow g)$. They receive the \emph{diagonals}
\begin{equation*}
\Delta_{f} := (\mathrm{id}_{X},f,\mathrm{id}_{f})\colon X \to (f \downarrow Y)
\ , \qquad \Delta_{g} := (g,\mathrm{id}_{Y},\mathrm{id}_{g})\colon Y \to (X
\downarrow g) \ ,
\end{equation*}
with
$p\Delta_{f} = \mathrm{id}_{X}$,
$q\Delta_{f} = f$,
$\alpha\Delta_{f} = \mathrm{id}_{f}$
and
$\bar p\Delta_{g} = g$,
$\bar q\Delta_{g} = \mathrm{id}_{Y}$,
$\beta\Delta_{g} = \mathrm{id}_{g}$.

\begin{lem}
\label{lem:comma-diagonal}
There are adjunctions
$\Delta_{f} \dashv p$
with identity unit and counit
$(\mathrm{id}_{p},\alpha)\colon \Delta_{f}p \Rightarrow \mathrm{id}$,
and
$\bar q \dashv \Delta_{g}$
with unit
$(\beta,\mathrm{id}_{\bar q})\colon \mathrm{id} \Rightarrow \Delta_{g}\bar q$
and identity counit. Hence restriction along $\Delta_{f}$ and $\Delta_{g}$ are
equivalences
\begin{equation*}
\mapp(sp,t) \xrightarrow{\ \simeq\ } \mapp(s,t\Delta_{f}) \ , \qquad
\mapp(t',s'\bar q) \xrightarrow{\ \simeq\ } \mapp(t'\Delta_{g},s') \ ,
\end{equation*}
for all
$s\colon X \to T$,
$t\colon (f \downarrow Y) \to T$,
$s'\colon Y \to T$
and
$t'\colon (X \downarrow g) \to T$.
\end{lem}

\begin{proof}
The pair
$(\mathrm{id}_{p},\alpha)$
is a $2$-morphism
$\Delta_{f}p = (p,fp,\mathrm{id}_{fp}) \Rightarrow (p,q,\alpha) = \mathrm{id}$.
It satisfies the triangle identities with the identity unit,
$p(\mathrm{id}_{p},\alpha) = \mathrm{id}_{p}$
and
$(\mathrm{id}_{p},\alpha)\Delta_{f} = (\mathrm{id}_{\mathrm{id}_{X}},
\mathrm{id}_{f}) = \mathrm{id}_{\Delta_{f}}$,
so it is the counit of an adjunction
$\Delta_{f} \dashv p$
by \cref{thm:adj-lifting}(2). By \eqref{eq:hom-adjunctions},
$p^{\natural} \dashv \Delta_{f}^{\natural}$
with identity unit, and the first equivalence is its transposition
$\theta \mapsto \theta\Delta_{f}$.
Likewise
$(\beta,\mathrm{id}_{\bar q})$
satisfies the triangle identities
$\bar q(\beta,\mathrm{id}_{\bar q}) = \mathrm{id}_{\bar q}$
and
$(\beta,\mathrm{id}_{\bar q})\Delta_{g} = \mathrm{id}_{\Delta_{g}}$
with the identity counit, so
$\bar q \dashv \Delta_{g}$
and
$\Delta_{g}^{\natural} \dashv \bar q^{\natural}$
with identity counit; the second equivalence is the inverse of its
transposition.
\end{proof}

\begin{lem}
\label{lem:comma-fibrewise}
\gutentag{001B}%
Let $F\colon A \to C$, $G\colon B \to C$,
$F'\colon A \to C'$ and $G'\colon B \to C'$ be functors of $\infty$-categories.
A functor
$(F \downarrow G) \to (F' \downarrow G')$
over $A \times B$ is an equivalence if and only if it induces equivalences
$\mapp(Fa,Gb) \to \mapp(F'a,G'b)$
on the fibres over all $(a,b)$.
\end{lem}

\begin{proof}
\gutentag{001C}%
$(F \downarrow G) \to A \times B$
is the pullback of
$\func(\Delta^{1},C) \to C \times C$
along $F \times G$, hence a bifibration \cite[Cor.~2.4.7.11]{htt},
with fibre $\mapp(Fa,Gb)$ over $(a,b)$. A functor over $A \times B$ between
bifibrations that is an equivalence on fibres is an equivalence
\cite[Prop.~2.4.7.6]{htt}.
\end{proof}

\begin{lem}
\label{lem:comma-maps}
Let
$h\colon X' \to Z$
and
$k\colon Y' \to Z$
be $1$-morphisms, and regard $(h \downarrow k)$, with structure $p'$, $q'$,
$\gamma\colon hp' \Rightarrow kq'$,
as a span from $X'$ to $Y'$.
\begin{enumerate}
\item Restriction along $\Delta_{f}$ is an equivalence over
$\Hom(X,X') \times \Hom(Y,Y')$
\begin{equation*}
\Hom_{\Span(\E)}\bigl((f \downarrow Y),(h \downarrow k)\bigr)
\xrightarrow{\ \simeq\ } (h_{\natural} \downarrow k_{\natural}f^{\natural}) \
.
\end{equation*}
Its inverse takes
$(a,b,\sigma\colon ha \Rightarrow kbf)$
to the map of spans
$c_{\sigma} := (ap,\,bq,\,kb\alpha \circ \sigma p)$.
\item Restriction along $\Delta_{g}$ is an equivalence over
$\Hom(X,X') \times \Hom(Y,Y')$
\begin{equation*}
\Hom_{\Span(\E)}\bigl((X \downarrow g),(h \downarrow k)\bigr)
\xrightarrow{\ \simeq\ } (h_{\natural}g^{\natural} \downarrow k_{\natural}) \
.
\end{equation*}
Its inverse takes
$(a,b,\tau\colon hag \Rightarrow kb)$
to the map of spans
$d_{\tau} := (a\bar p,\,b\bar q,\,\tau\bar q \circ ha\beta)$.
\end{enumerate}
\end{lem}

\begin{proof}
(1) By the universal property of $(h \downarrow k)$, a map of spans
$(f \downarrow Y) \to (h \downarrow k)$
over $(a,b)$ is a $1$-morphism
$(ap,bq,\theta)$
with
$\theta\colon hap \Rightarrow kbq$.
So the source is the comma $\infty$-category
$(h_{\natural}p^{\natural} \downarrow k_{\natural}q^{\natural})$
of triples $(a,b,\theta)$, and restriction along $\Delta_{f}$ sends
$(a,b,\theta)$
to
$(a,b,\theta\Delta_{f}\colon ha \Rightarrow kbf)$,
as
$p\Delta_{f} = \mathrm{id}_{X}$
and
$q\Delta_{f} = f$.
On the fibres over $(a,b)$ this is
$\mapp(hap,kbq) \to \mapp(ha,kbf)$,
$\theta \mapsto \theta\Delta_{f}$,
an equivalence by \cref{lem:comma-diagonal} with $s = ha$ and $t = kbq$. By
\cref{lem:comma-fibrewise}, restriction along $\Delta_{f}$ is an equivalence.
It takes $c_{\sigma}$ to
$(a,bf,kb\alpha\Delta_{f} \circ \sigma) = (a,bf,\sigma)$,
as
$\alpha\Delta_{f} = \mathrm{id}_{f}$.

(2) Similarly, a map of spans
$(X \downarrow g) \to (h \downarrow k)$
over $(a,b)$ is a $1$-morphism
$(a\bar p,b\bar q,\theta)$
with
$\theta\colon ha\bar p \Rightarrow kb\bar q$,
and restriction along $\Delta_{g}$ sends it to
$(a,b,\theta\Delta_{g}\colon hag \Rightarrow kb)$.
On fibres this is
$\mapp(ha\bar p,kb\bar q) \to \mapp(hag,kb)$,
an equivalence by \cref{lem:comma-diagonal} with $t' = ha\bar p$ and
$s' = kb$, so restriction along $\Delta_{g}$ is an equivalence by
\cref{lem:comma-fibrewise}. It takes $d_{\tau}$ to
$(ag,b,\tau \circ ha\beta\Delta_{g}) = (ag,b,\tau)$,
as
$\beta\Delta_{g} = \mathrm{id}_{g}$.
\end{proof}

\begin{defi}
\label{def:transposing-adjunction}
\gutentag{001D}%
Let $\E$ have lax pullbacks.
\begin{enumerate}
\item A \emph{transposing adjunction} between $f\colon X \to Y$ and $g\colon Y
\to X$ is an equivalence
$e\colon (f \downarrow Y) \xrightarrow{\simeq} (X \downarrow g)$
in $\Span_{\E}(X,Y)$.
\item Let $\TAdj(\E)$ be the full sub-$(\infty,2)$-category of
$\func(\Delta^{1},\Span(\E))$
on the transposing adjunctions.
\end{enumerate}
\end{defi}

\begin{lem}
\label{lem:tadj-homs}
\gutentag{0100}%
Let $e$ and $e'$ be transposing adjunctions between
$f\colon X \to Y$ and $g\colon Y \to X$ and between
$f'\colon X' \to Y'$ and $g'\colon Y' \to X'$.
A $1$-morphism $(c,d)\colon e \to e'$ of $\TAdj(\E)$ consists of maps of spans
$c$ and $d$ over some
$(a,b) \in \Hom(X,X') \times \Hom(Y,Y')$
and an equivalence $de \simeq e'c$:
\begin{equation*}
\begin{tikzcd}
(f \downarrow Y) \ar[r, "e"] \ar[d, "c"'] & (X \downarrow g) \ar[d, "d"] \\
(f' \downarrow Y') \ar[r, "e'"'] & (X' \downarrow g')
\end{tikzcd}
\end{equation*}
Its hom-$\infty$-category is the pullback
\begin{equation*}
\Hom_{\TAdj(\E)}(e,e') \simeq \Hom_{\Span(\E)}\bigl((f \downarrow Y),(f'
\downarrow Y')\bigr) \times_{\Hom_{\Span(\E)}((f \downarrow Y),(X' \downarrow
g'))} \Hom_{\Span(\E)}\bigl((X \downarrow g),(X' \downarrow g')\bigr)
\end{equation*}
along $e'_{\natural}$ and $e^{\natural}$, over
$\Hom(X,X') \times \Hom(Y,Y')$.
\end{lem}

\begin{proof}
\gutentag{0101}%
Apply \cite[Prop.~4.1.2]{abellangagnahaugseng2024straightening} to
$\func(\Delta^{1},\Span(\E))$,
as for $\Span(\E)$ above.
\end{proof}

\gutentag{001E}%
For $\E = \Cat$, $(f \downarrow Y)$ and $(X \downarrow g)$ are bifibrations over
$X \times Y$ with fibres $\mapp_{Y}(fx,y)$ and $\mapp_{X}(x,gy)$,
so by \cref{lem:comma-fibrewise} a transposing adjunction is a functor over $X
\times Y$ inducing equivalences
$\mapp_{Y}(fx,y) \simeq \mapp_{X}(x,gy)$.
The name is taken from Riehl--Shulman's synthetic setting
\cite[Def.~11.1]{riehlshulman2017synthetic}.

\begin{thm}[adjunctions via lax pullbacks]
\label{thm:comma-adjunctions}
\gutentag{001F}%
Let $\E$ be an $(\infty,2)$-category with lax pullbacks, and let $f\colon X \to
Y$ and $g\colon Y \to X$. Every $1$-morphism
$(f \downarrow Y) \to (X \downarrow g)$
of $\Span_{\E}(X,Y)$ is, for a unique
$\eta\colon \mathrm{id}_{X} \Rightarrow gf$,
the map
\begin{equation*}
e_{\eta}\colon (f \downarrow Y) \to (X \downarrow g) \ ,
\qquad (u,v,\varphi\colon fu \Rightarrow v) \mapsto (u,v,g\varphi \circ \eta
u\colon u \Rightarrow gv) \ ;
\end{equation*}
that is,
$\eta \mapsto e_{\eta}$
is an equivalence
$\mapp(\mathrm{id}_{X},gf) \simeq \Hom_{\Span_{\E}(X,Y)}((f \downarrow Y),(X
\downarrow g))$.
The map $e_{\eta}$ is a transposing adjunction if and only if $\eta$ is the unit
of an adjunction $f \dashv g$, and
$A \mapsto e_{A(\eta)}$
is an equivalence from the space of adjunctions $A$ with $A(l) = f$ and $A(r) =
g$ to the space of transposing adjunctions between $f$ and $g$.
In particular, $f$ is a left adjoint if and only if there is a transposing
adjunction between $f$ and some $g$, and $g$ is then a right adjoint of $f$.
\end{thm}

\gutentag{0020}%
So when $\E$ has lax pullbacks, an adjunction can be defined as a transposing
adjunction, without reference to $\Adj$. For $\E = \Cat$, $e_{\eta}$ is
$\varphi \mapsto g\varphi \circ \eta_{x}$
on the fibres, so \cref{thm:comma-adjunctions} says that $\eta$ is the unit of
an adjunction if and only
if it is a unit transformation in the sense of \cite[Def.~5.2.2.7]{htt}.

\begin{proof}[Proof of \cref{thm:comma-adjunctions}]
\gutentag{0021}%
\emph{Step 1: the maps between $(f \downarrow Y)$ and $(X \downarrow g)$.}
The fibre of \cref{lem:comma-maps}(1) for $h = \mathrm{id}_{X}$ and $k = g$,
i.e.\
$(h \downarrow k) = (X \downarrow g)$,
over
$(a,b) = (\mathrm{id}_{X},\mathrm{id}_{Y})$
is an equivalence
\begin{equation*}
\Hom_{\Span_{\E}(X,Y)}\bigl((f \downarrow Y),(X \downarrow g)\bigr)
\xrightarrow{\ \simeq\ } \mapp(\mathrm{id}_{X},gf) \ .
\end{equation*}
It takes a map of spans $e$ to the $2$-morphism $\eta$ with
$e\Delta_{f} = (\mathrm{id}_{X},f,\eta)$,
and its inverse takes $\eta$ to
$c_{\eta} = (p,q,g\alpha \circ \eta p)$,
which is the map $e_{\eta}$ of the statement. This is the first claim.
Likewise, the fibre of \cref{lem:comma-maps}(2) for $h = f$ and
$k = \mathrm{id}_{Y}$
over
$(\mathrm{id}_{X},\mathrm{id}_{Y})$
shows that the maps of spans
$(X \downarrow g) \to (f \downarrow Y)$
are the
\begin{equation*}
d_{\varepsilon} = (\bar p,\bar q,\varepsilon\bar q \circ f\beta)\colon
(u,v,\psi) \mapsto (u,\,v,\,\varepsilon v \circ f\psi) \ , \qquad
d_{\varepsilon}\Delta_{g} = (g,\mathrm{id}_{Y},\varepsilon) \ ,
\end{equation*}
for
$\varepsilon\colon fg \Rightarrow \mathrm{id}_{Y}$.
By \cref{lem:comma-maps}(1) for $h = f$ and $k = \mathrm{id}_{Y}$, an
endomorphism $c$ of $(f \downarrow Y)$ in $\Span_{\E}(X,Y)$ is equivalent to
the identity if and only if
$c\Delta_{f} \simeq \Delta_{f}$,
and by \cref{lem:comma-maps}(2) for $h = \mathrm{id}_{X}$ and $k = g$, an
endomorphism $d$ of $(X \downarrow g)$ is equivalent to the identity if and only
if
$d\Delta_{g} \simeq \Delta_{g}$.

\emph{Step 2: $e_{\eta}$ is an equivalence if and only if $\eta$ is a unit.} As
equivalences in $\Span_{\E}(X,Y)$ are detected in $\E$,
$e_{\eta}$ is an equivalence if and only if it has an inverse in
$\Span_{\E}(X,Y)$, which by Step 1 is some $d_{\varepsilon}$. By Step 1,
$d_{\varepsilon}$ is inverse to $e_{\eta}$ if and only if
\begin{equation*}
d_{\varepsilon}e_{\eta}\Delta_{f} = (\mathrm{id}_{X},\,f,\,\varepsilon f \circ
f\eta) \simeq \Delta_{f} \ , \qquad e_{\eta}d_{\varepsilon}\Delta_{g} =
(g,\,\mathrm{id}_{Y},\,g\varepsilon \circ \eta g) \simeq \Delta_{g} \ ,
\end{equation*}
i.e.\ if and only if $\eta$ and $\varepsilon$ satisfy the triangle identities.
So $e_{\eta}$ is an equivalence if and only if $\eta$ is the unit of an
adjunction $f \dashv g$ (\cref{thm:adj-lifting}(2)), and then
$e_{\eta}^{-1} = d_{\varepsilon}$
for its counit $\varepsilon$.

\emph{Step 3: adjunctions are transposing adjunctions.} Let $\mathrm{Ad}(f)$ be
the fibre of $\ell^{*}$ over $f$, the space of adjunctions $A$ with $A(l) = f$,
and let $\mathrm{Un}(f)$ be the space of pairs $(g,\eta)$ with $\eta$ the unit
of an adjunction $f \dashv g$. Both lie over $\mapp(Y,X)$, and
$A \mapsto (A(r),A(\eta))$
is a map
$\mathrm{Ad}(f) \to \mathrm{Un}(f)$
over $\mapp(Y,X)$. By Steps 1 and 2,
$\eta \mapsto e_{\eta}$
identifies the fibre of $\mathrm{Un}(f)$ over $g$ with the space of transposing
adjunctions between $f$ and $g$, so it suffices to show that
$\mathrm{Ad}(f) \to \mathrm{Un}(f)$
is an equivalence, which is \cref{thm:adj-lifting}(3).
\end{proof}

\begin{rmk}
\label{rmk:comma-literature}
\gutentag{0022}%
Riehl--Shulman compare transposing adjunctions with adjunctions with a given
unit in their synthetic setting \cite[Cor.~11.21,
Thm.~11.23]{riehlshulman2017synthetic}. When $\E$ is presented by an
$\infty$-cosmos, \cref{thm:comma-adjunctions} is
\cite[Prop.~4.1.1]{riehlverity2022elements},
with the strict comma objects of the cosmos;
for $\E = \Cat$ see \cite[Prop.~4.4.3]{riehlverity2015twocategory}.
If $\E$ does not have lax pullbacks, one can apply
\cref{thm:comma-adjunctions} to $\Hom_{\E}(-,f)$ and $\Hom_{\E}(-,g)$ in
$\func(\E\opcat,\Cat)$, into which $\E$ embeds fully faithfully by the enriched
Yoneda lemma \cite{hinich2020yoneda,heine2023enriched}.
\end{rmk}
